\section{Introduction}\label{sec:intro}

Cartel enforcement is designed to restore competition. But enforcement has side effects. When a convicted firm contracts---losing contracts, paying fines, facing reputational damage---its workers leave. In thin labor markets, those workers have few alternative employers, and they flow disproportionately to the cartel firm's competitors in the same product market. If these workers carry coordination know-how---pricing strategies, bid-rotation protocols, information about rivals' costs---the receiving firms may become more collusive, not more competitive. Enforcement would then contain an ironic seed: dismantling one cartel while planting the next.

This concern is not hypothetical. \citet{lamichhane2024collusion} provide experimental evidence that worker mobility between firms facilitates collusive pricing in Bertrand markets: workers carry pricing-strategy information, and firms that exchange employees coordinate on higher prices and more market sharing. \citet{herrera2025collusion} find that common leadership---shared directors and executives---raises the probability of product-market collusion by 11~percentage points. Both papers establish that inter-firm personnel connections are a channel for collusion. What neither paper can test is whether \emph{enforcement-driven} labor reallocation activates this channel in practice. That is the question we address.

We link seven procurement cartels convicted by Brazil's antitrust authority (CADE) to 2.6~million competitive auctions on S\~{a}o Paulo's electronic procurement platform (BEC) and to the universe of matched employer--employee records (RAIS). When cartel conduct periods end, 34{,}914 workers flow from the convicted firms to 3{,}960 other BEC-active firms. Of these receiving firms, 962 bid on the same item codes as the cartel firms both before and after absorbing ex-cartel workers, creating a natural experiment: does the arrival of ex-cartel workers change the receiving firm's competitive behavior?

The answer is yes---but not in the way the contagion hypothesis predicts. Using a stacked difference-in-differences design with never-treated controls, we find that receiving firms increase their win rates by 5--6~percentage points in the product markets where the cartel operated. All five post-treatment coefficients are significant at the 1\% level, and the effect is growing over time. But receiving firms do not raise their prices: the log price ratio relative to the reference price is unchanged after absorbing ex-cartel workers. The apparent price increase in a standard two-way fixed-effects specification ($+0.26$ log points, highly significant) is an artifact of staggered-treatment bias; correcting for this via the stacked estimator eliminates the effect entirely ($+0.05$, insignificant).

Three additional tests pin down the mechanism. \emph{First}, an exit placebo compares firms that absorb workers from non-cartel firms that simply stop bidding on BEC. The win-rate effect is 3.7 times larger for cartel-worker absorbers than for exit-worker absorbers ($+3.2$~pp vs.\ $+0.8$~pp in CEM-matched samples), confirming that ex-cartel workers bring something that generic labor-market movers do not. The price effect, by contrast, is identical for both groups---any firm that absorbs workers from an exiting competitor charges more, regardless of cartel status. \emph{Second}, characterizing the movers reveals that managers and professionals are \emph{under}represented among those who leave cartel firms ($9.9\%$) relative to those who stay ($12.9\%$). The modal mover is a service or sales worker earning 2.4 minimum wages. This is a displacement profile, not a knowledge-transfer profile. \emph{Third}, CEM matching on pre-period characteristics (bid volume, item breadth, win rate, firm size) eliminates a pre-trend at $\tau = -3$ in the unmatched specification, confirming that the treatment effect is not driven by selection on trajectory.

The combined evidence points to \textbf{capacity transfer}, not collusion contagion. Ex-cartel workers bring operational know-how---familiarity with specific products, supply chains, and procurement procedures---that helps the receiving firm compete more effectively in the markets the cartel once dominated. They do not bring the coordination capability that would allow the receiving firm to charge supra-competitive prices. Enforcement dismantles the cartel and, through labor reallocation, strengthens the cartel's former rivals---without transmitting the ability to collude.

The paper contributes to three literatures. To the \textbf{cartel enforcement} literature, we provide the first evidence on whether enforcement-driven labor reallocation creates collusion spillovers. The answer is no: the reallocation is competitive, not collusive. This is a policy-relevant null: competition authorities need not worry that debarment or conviction seeds the next cartel through worker mobility. To the \textbf{worker mobility and knowledge transfer} literature \citep{stoyanov2012worker, serafinelli2019good, jarosch2021learning}, we show that the type of knowledge that transfers depends on which workers move. When operational staff move, the effect is capacity deepening; when managers move (which happens less), the channel is attenuated. To the \textbf{econometric} literature on staggered difference-in-differences \citep{goodman2021difference, cengiz2019effect}, we provide a vivid example of how TWFE bias can generate a false positive: the price ``contagion'' that appears in the naive specification ($+0.26$, $t > 5$) vanishes entirely under the stacked estimator ($+0.05$, insignificant).

The remainder of the paper proceeds as follows. Section~\ref{sec:background} describes the institutional setting. Section~\ref{sec:data} describes the data and sample construction. Section~\ref{sec:strategy} presents the empirical strategy. Section~\ref{sec:results} reports results. Section~\ref{sec:conclusion} concludes.
